\chapter{Literature Review}

In the developmental history of stochastic volatility models, the 3/2 model has gradually become a research focal point in the fields of option pricing and risk management within financial engineering, owing to its unique nonlinear diffusion structure and superior ability to capture extreme market volatility. This chapter follows a logical thread of ``model paradigm evolution — theoretical framework development — pricing method iteration — market application expansion — domestic research status.'' It systematically reviews the domestic and international research achievements related to the 3/2 model, ultimately clarifying the core gaps in existing research through a literature summary to establish the academic positioning and marginal contributions of this study.

\section{Evolution of Stochastic Volatility Models}

The evolution of option pricing theory is essentially a process of market participants continuously deepening their understanding of the ``non-normalit'' of financial asset returns and the ``time-varying risk'' of volatility. Since the Black-Scholes-Merton (BSM) model \citet{blackPricingOptionsCorporate1973,mertonTheoryRationalOption1973} established the basic paradigm of arbitrage-free pricing, its core assumptions—constant volatility and the underlying price following geometric Brownian motion—have provided a concise and efficient analytical framework for derivatives pricing, laying the foundation of modern financial engineering. However, as the complexity of global financial markets continues to escalate, the core assumptions of the BSM model have systematically diverged from true market characteristics. A large number of empirical studies have shown that the return distributions of various global assets universally exhibit significant ``heavy-tailed'' characteristics. Meanwhile, implied volatility in the options market presents a ``volatility smile'' varying with strike prices and a ``volatility skew'' varying with time to maturity \citet{rubinsteinNonparametricTestsAlternative1985}. These phenomena are particularly prominent during market crises and extreme conditions, revealing the nonlinear pricing characteristics of the market's tail risk premium and driving a paradigm shift in volatility modeling from determinism to stochasticity.

To correct the theoretical deficiencies of the BSM model, academia first developed local volatility models. \citet{dermanRidingSmile1994} and \citet{Dupire1994} constructed deterministic local volatility frameworks dependent on the underlying price and time from the perspectives of discrete trees and continuous partial differential equations, respectively, which can perfectly fit the implied volatility surface at a specific moment. However, local volatility models fundamentally remain within the realm of deterministic volatility; they cannot characterize the stochastic fluctuation features of volatility itself and possess weak predictive power for future volatility surfaces, making it difficult to meet the practical demands of dynamic hedging and long-term derivatives pricing.

Against this backdrop, Stochastic Volatility (SV) models emerged. Their core concept is to define volatility itself as a stochastic diffusion process, thereby simultaneously characterizing the stochastic fluctuations of the underlying price and the time-varying features of volatility itself. Affine SV Models, represented by \citet{hestonSimpleNewFormula1997}, introduced a square-root diffusion process (also known as the 1/2 model) to characterize the mean-reverting property of volatility, deriving an affine closed-form solution for the characteristic function of the underlying log-price. This greatly reduced the computational complexity of model calibration and option pricing, becoming a standard pricing cornerstone in the global derivatives industry. Subsequently, targeting the complex volatility dynamics in interest rate and foreign exchange markets, the SABR model proposed by \citet{Hagan2002} further enhanced the ability to capture the shape of implied volatility curves by introducing a stochastic volatility term and a correlation structure between volatility and the underlying price, becoming the mainstream model for OTC interest rate derivatives pricing.

However, with the evolution of market microstructure and the frequent occurrence of extreme market conditions, the endogenous limitations of the traditional 1/2 model have gradually become apparent. Scholars have found that the Heston model possesses theoretical shortcomings in fitting the steep implied volatility of ultra-short-term options, and its linear drift term struggles to explain the explosive growth and nonlinear mean-reverting characteristics of volatility during extreme crises \citep{jonesDynamicsStochasticVolatility2003}. This drove academia to transition from affine frameworks to non-affine frameworks. As a typical representative of non-affine SV models, the 3/2 model features a diffusion term proportional to the instantaneous variance raised to the power of 3/2, and a drift term exhibiting a nonlinear mean-reverting structure. This stronger nonlinearity endows the model with a more flexible kurtosis characterization and tail risk capture capability. Consequently, when dealing with scenarios involving severe underlying price jumps and highly clustered volatility, it demonstrates a deeper theoretical foundation and superior empirical performance compared to the traditional 1/2 model.

\section{Development of the 3/2 Model and Its Extended Frameworks}

Since \citet{lewisOptionValuationStochastic2000} systematically derived the fundamental analytical properties of the 3/2 stochastic volatility model, academic research has extensively explored the theoretical framework optimization, analytical property expansion, and numerical implementation schemes of the 3/2 model, forming a relatively complete system of theoretical achievements.

In terms of constructing the model's fundamental analytical framework, \citet{carrNewApproachOption2007} developed a novel pricing framework by modeling the variance swap rate directly. This bypassed the complex process of modeling the market price of instantaneous volatility risk, providing a completely new implementation path for the risk-neutral pricing of the 3/2 model. Addressing the option pricing and calibration challenges brought by the 3/2 model's non-affine nature, \citet{medvedevApproximationCalibrationShortTerm2007} proposed a comprehensive analytical framework based on short-maturity asymptotic expansions. They derived a closed-form expansion for implied volatility under non-affine local volatility models, providing a core theoretical tool for the rapid cross-sectional calibration of non-affine models, including the 3/2 model, and serving as a crucial foundation for subsequent empirical research on the 3/2 model.

Regarding the extension and optimization of the model structure, to achieve consistent modeling of spot equities and VIX derivatives, \citet{baldeauxConsistentModelingVIX2012} introduced a simultaneous jump term into the pure diffusion 3/2 model, constructing the ``3/2 plus jumps model'' to further enhance the model's ability to characterize extreme volatility in severe markets. \citet{grasselli42Stochastic2017}, in their pioneering research, proposed the ``4/2 stochastic volatility model,'' which defines the diffusion term of the instantaneous variance as a linear combination of a 1/2 term and a 3/2 term. This organically integrated the Heston model and the 3/2 model within the same framework, forming a generalized SV model framework that effectively combines the advantages of both models, greatly expanding the theoretical application boundaries of the 3/2 model. To tackle the simulation difficulties caused by the 3/2 model's non-affine characteristics, \citet{kouarfateExplicitSolutionSimulation2021} developed a novel simulation framework based on an explicit solution for the 3/2 model's variance process. This framework ingeniously leverages the core mathematical property that the variance process of the 3/2 model is the reciprocal of a CIR process, providing a new implementation approach for Monte Carlo simulations of the 3/2 model.

\section{Evolution of Exact Simulation Methods}

In the pricing practice of stochastic volatility models, Monte Carlo simulation is the core tool for handling path-dependent derivatives and high-dimensional model pricing. The central developmental direction of simulation methods is to eliminate time-discretization errors while balancing computational efficiency and full-scenario numerical stability. Since exact simulation methods were introduced, academia has conducted continuous optimization research on exact simulation schemes for both affine models and the non-affine 3/2 model.

In foundational research on the exact simulation of affine SV models, \citet{broadieExactSimulationStochastic2006} first proposed an unbiased exact simulation method for the underlying stock price and variance process under the Heston model and other affine jump-diffusion processes, known as the Exact MC scheme. By sampling the conditional distribution of the integrated variance, this method completely eliminated the systematic errors caused by time discretization, laying the theoretical foundation for the exact simulation of SV models. Subsequently, \citet{glassermanGammaExpansionHeston2011} significantly optimized the Broadie-Kaya algorithm. By utilizing \citet{pitmanDecompositionBesselBridges1982}'s conclusions on the decomposition of integral series of diffusion processes, they decomposed the conditional integrated variance into an infinite series of Gamma random variables. This avoided the computationally expensive numerical inverse Laplace transform in the original algorithm, markedly improving the computational efficiency of exact simulations.

Focusing on the exact simulation of the non-affine 3/2 model, \citet{baldeauxEXACTSIMULATION322012} conducted pioneering research, exploring the exact simulation scheme for the stock price process under the 3/2 model for the first time. Utilizing conclusions regarding the transition density of diffusion processes derived via Lie symmetry analysis by \citet{craddockCalculationExpectationsClasses2009}, this study successfully adapted and applied \citet{broadieExactSimulationStochastic2006}'s exact simulation algorithm for affine processes to the 3/2 model. They derived a closed-form solution for the conditional Laplace transform of the integrated variance under the 3/2 model, providing the first theoretically viable implementation scheme for its unbiased simulation.

In recent years, exact simulation algorithms for the 3/2 model have achieved further breakthroughs in numerical stability and computational efficiency. \citet{choiSimulationSchemesHeston2024} proposed a novel exact simulation scheme that does not require the use of highly oscillatory modified Bessel functions. Their core finding was that when conditioning on the Poisson variable used to simulate the terminal variance, the conditional integrated variance of the 3/2 model can be significantly simplified, fundamentally resolving the core pain point of modified Bessel functions easily causing numerical overflow under extreme parameters in the original Baldeaux algorithm. Additionally, \citet{brignoneExactSimulationStochastic2026} proposed a new exact simulation framework utilizing the conditional Fourier-Cosine (COS) method. By constructing the conditional cumulative distribution function of the integrated variance through cosine series expansion, they achieved highly efficient and accurate sampling of the integrated variance. This resolved the long-standing implementation difficulties of error control and high computational overhead in exact simulations since \citet{broadieExactSimulationStochastic2006}, further expanding the engineering application boundaries of exact simulation methods.

\section{Research on Fourier Transform Methods}

The Fourier transform pricing method is the core technical pathway for the rapid pricing and comprehensive calibration of European options under SV models. Its central idea is to convert the risk-neutral expectation of the option price into a complex-domain integral of the underlying log-price's characteristic function, achieving rapid solutions through numerical integration algorithms and completely circumventing the statistical and time-discretization errors of Monte Carlo simulations.

In foundational research on Fourier transform pricing methods, \citet{carrOptionValuationUsing1999} proposed an option pricing method based on the Fast Fourier Transform (FFT). By introducing a damping factor to resolve the non-integrability issue of the option payoff function's Fourier transform, they derived a closed-form inverse Fourier transform expression for European option prices. Combined with the FFT algorithm, this enabled the simultaneous calculation of option prices across a full range of strike prices within an $O(N \log N)$ time complexity, becoming the mainstream engineering method for option pricing under SV models. Addressing the limitations of the FFT method, such as the picket-fence effect and the restriction to equally spaced strike prices, \citet{fangNovelPricingMethod2009} proposed the COS pricing method based on Fourier cosine series expansions. By expanding the probability density function of underlying returns into a cosine series over an effective support interval, they derived a closed-form series summation for option prices. This achieved exponential convergence rates while allowing flexible calculations for arbitrary strike prices, demonstrating extremely strong applicability and computational efficiency advantages in non-affine SV models.

For the 3/2 model, due to its non-affine structural nature, the underlying log-price lacks the simple exponential-affine characteristic function found in affine models. The analytical expression of its characteristic function involves complex special functions, posing certain challenges to the application of Fourier transform methods. \citet{lewisOptionValuationStochastic2000}'s early research on complex-domain integral pricing provided the core theoretical framework for exploring the analytical properties of the 3/2 model; \citet{dufresne2001integrated}'s research further pointed out that the integral distribution of the 3/2 process is closely related to Whittaker functions and modified Bessel functions. He fully derived the distributional properties of the integrated variance under the 3/2 model, laying the mathematical foundation for the closed-form solution of the characteristic function. In practical applications, the FFT pricing method is widely used for parameter calibration and option pricing of the 3/2 model, but it still faces challenges of insufficient numerical stability and excessive computation time under extremely short maturities and extreme volatility parameters \citet{drimusOptionsRealizedVariance2011}. In contrast, the COS method, with its exponential convergence rate and strong numerical stability, has demonstrated significant performance advantages in the pricing and calibration of the 3/2 model, becoming a primary focus of research in this field in recent years.

\section{Research on Stochastic Volatility Models in the Cryptocurrency Options Market}

As crypto assets represented by Bitcoin (BTC) and Ethereum (ETH) enter the era of institutional trading, the scale of the cryptocurrency options market has grown exponentially, and its unique market dynamics have become a frontier research direction in financial engineering. Unlike traditional equity and foreign exchange markets, the cryptocurrency market is characterized by $7\times24$ continuous trading, no price limits, extremely high endogenous volatility, and frequent extreme market events. Its return distribution exhibits significant heavy-tailed features and an extremely high volatility of volatility (Vol-of-Vol, VoV), posing severe challenges to traditional option pricing models.

\citet{houPricingCryptocurrencyOptions2020}'s research points out that the volatility of crypto assets exhibits stronger persistence and a higher Vol-of-Vol than traditional stock assets, with extreme tail risks significantly higher than those in traditional financial markets. This frequently causes severe valuation deviations in crypto option pricing when using the BSM model or even the basic Heston model; especially under extreme conditions, the models' ability to fit the volatility smile sharply declines. Addressing this phenomenon, academia primarily explores model frameworks suitable for the crypto market from two directions: one stream of research attempts to capture the instantaneous extreme jumps in cryptocurrency prices by introducing jump-diffusion processes into SV models \citet{scaillet2020high}; the other focuses on improving the nonlinear structure of the volatility process, exploring the adaptability of non-affine SV models in the crypto market.

In recent years, due to the 3/2 model's natural advantages in characterizing implied volatility skew and capturing nonlinear mean-reverting features, research on its adaptability to the cryptocurrency market has significantly increased. Studies show that when crypto assets face regulatory policy shifts or macro liquidity shocks, the mean reversion of their volatility exhibits obvious nonlinear acceleration, which aligns highly with the mathematical properties of the 3/2 model (Fassas et al., 2022). Furthermore, given the crypto market's unique $7\times24$ continuous trading mechanism, utilizing the 3/2 model combined with high-frequency data for real-time volatility forecasting and dynamic option hedging has become a current research hotspot in the field. However, overall, existing research on the 3/2 model in the crypto market mostly focuses on static parameter calibration and model fit comparisons, lacking dynamic evolutionary analysis of model parameters throughout the full cycle of extreme crash events. Moreover, there has been no systematic empirical testing of the performance of different pricing methods under extreme market conditions. This research gap provides critical space for the empirical analysis conducted in this thesis.

\section{Literature Review Summary and Research Positioning of this Paper}

In summary, international academia has conducted extensive research around the theoretical construction, pricing method iteration, and market application expansion of the 3/2 stochastic volatility model. Significant achievements have been made in exploring analytical properties, constructing extended frameworks, optimizing exact simulation methods, and improving Fourier pricing methods, laying a solid theoretical foundation for this study. Concurrently, domestic research largely focuses on traditional SV models like the Heston and SABR models, while systematic research on the 3/2 model remains relatively scarce. There are still prominent research gaps concerning performance evaluations of pricing methods in extreme market environments and empirical adaptability tests in high-volatility crypto asset markets.

Despite the relative abundance of existing research, three core shortcomings remain. First, performance evaluations of existing 3/2 model pricing methods are largely confined to idealized cases using benchmark parameters from literature. There is a lack of full-scenario numerical stability testing covering both conventional markets and extreme parameter combinations. Some methods experience issues like numerical overflow and computational failure under real-world extreme parameters, such as extreme Vol-of-Vol, strong correlation coefficients, and extremely short maturities, leaving their engineering applicability boundaries vaguely defined. Second, existing research has not yet conducted empirical re-evaluations of pricing methods using parameters calibrated from genuinely extreme markets. Performance in theoretical cases cannot fully equate to practical efficacy in real markets, and empirical evidence is still lacking regarding the differences in accuracy, efficiency, and robustness among various pricing methods during extreme conditions. Third, empirical research on the 3/2 model in highly volatile cryptocurrency markets remains imperfect, lacking dynamic evolutionary analysis of model parameters and volatility smile fitting tests over the full cycle of extreme crash events; the pure diffusion 3/2 model's adaptability to extreme markets has not yet been fully validated.

Based on the deficiencies in existing research, this paper clarifies its core research positioning: taking European option pricing under the 3/2 stochastic volatility model as the research subject, this study systematically constructs a unified implementation framework for eleven mainstream pricing methods across three paradigms. Through multi-scenario numerical experiments and empirical testing in extreme markets, it comprehensively evaluates the numerical stability, pricing accuracy, and computational efficiency of each method to clarify their engineering applicability boundaries. Concurrently, using a severe Bitcoin crash event as a sample, it validates the 3/2 model's adaptability to extreme markets and identifies the time-varying patterns of its model parameters. This research aims to bridge the gaps in existing studies regarding performance testing in extreme scenarios and empirical analysis of real markets, providing reliable method selection and empirical support for 3/2 model pricing practices in high-volatility markets.